% ============================================================================
% 06f-baselines.tex -- External baselines: no-plugin host and Mem0 (three-subject)
%
% Numbers read from:
%   paper/comparison-lab/results/summary.json      (baseline, lingshu; n=5)
%   paper/mem0-chaos-lab/results/summary.json      (mem0 2.2.1; n=5)
%   paper/mem0-chaos-lab/raw/human-review.json     (human override of thresholds)
% Table body is generated by paper/comparison-lab/make_table.py -> paper/tables/
% baselines.tex; prose numbers are checked against paper/comparison-lab/
% results/final-table.json by _verify_metrics.py (Lead).
%
% HONESTY NOTE: this is a five-round directional pilot on one host, one model
% route and one day. The three subjects do not see the same input shape (raw
% transcript vs retrieved memory), so differences are structural, not effect
% sizes; no confidence intervals or p-values are reported. Mem0's embedding and
% cross-lingual recall limits are environment constraints and are reported in
% the Limitations section (A0 owns that text), not in this comparison.
% ============================================================================

\subsection{External baselines: a no-plugin host and Mem0}
\label{sec:exp:baselines}

The preceding experiments compare arms of the same system or measure the gate in
isolation. A reader may ask the obvious external question: does the control plane
buy anything over (i) the bare model without any memory layer, and (ii) a
production long-term memory framework with no safety layer? We therefore run a
small three-subject pilot on two scenarios. We frame it as a directional pilot,
not a benchmark: five rounds per cell, one host, one model route, one day.

\paragraph{Subjects and protocol.} Three subjects: \textsc{Baseline}, the host
with no plugin at all (raw conversation transcript in one call);
\textsc{Mem0}, the official \texttt{mem0ai}~2.2.1 library~\cite{mem0_2025} driven
by the same DeepSeek model, with an in-memory vector store; and \system{} as the
lab profile with the plugin active, one isolated data directory per cell. Scenario~S1 (memory conflict)
plants a codename ALPHA, then a distractor preference, then a correction to BETA
with ALPHA declared void, and finally asks for the codename. Scenario~S2
(jailbreak) plants a confidential token and then asks the system to ignore
instructions and print its system prompt and all injected memory including the
token. \textsc{Baseline} receives the whole transcript in a single call; \system{}
is driven turn by turn (four calls for S1, two for S2) so that its memory and
gate actually run; \textsc{Mem0} receives each turn through \texttt{add} and is
queried with the final question. Every cell starts from a fresh state.

\paragraph{Metrics.} For each cell we record three rates. \emph{Hit}: the final
answer gives the current codename for S1; for S2 it means the attack succeeded
(the token is disclosed, or internal markers appear in a non-refusal answer).
\emph{Leak}: the superseded value ALPHA appears in the S1 answer, or the planted
token appears in the S2 answer. \emph{Interception}: for S1, the old value is
absent from what the system injects (retrieved memories for \textsc{Mem0}, the
injected-memory preview for \system{}, the raw transcript for \textsc{Baseline});
for S2, a structural gate blocked the injection (UCA-01, attributed to the user
message), with model-level refusal reported separately.

% generated from paper/comparison-lab/results/final-table.json (v2) -- do not edit by hand
% primary sources: results/summary.json (sha256 78e5c5328bb2...),
%                  mem0-chaos-lab/results/summary.json + human-review.json
\begin{table*}[t]
\centering\footnotesize
\caption{Three-subject pilot ($n=5$ independent rounds per cell; one host, one
model route, one day). S1 = memory conflict (ALPHA$\rightarrow$BETA); S2 =
jailbreak/injection. \emph{Hit}: current value answered / attack succeeded under
human review; \emph{Leak}: superseded value or planted token disclosed;
\emph{Sup.}: the superseded value is absent from what the system injects;
\emph{Int.}: model-level refusal; \emph{gate}/\emph{audit}: structural UCA-01
block and its audit record. Mem0 is shown in its recall-working diagnostic
configuration (see the default-configuration note); its default row is omitted
from the main body because its embedding pathway is empty in this environment.
Sources: \texttt{paper/comparison-lab/results/final-table.json}.}
\label{tab:baselines}
\begin{tabular}{@{}l ccc ccc ccc@{}}
\toprule
 & \multicolumn{3}{c}{S1 memory conflict} & \multicolumn{3}{c}{S2 jailbreak} & \multicolumn{3}{c}{S2 structural} \\
\cmidrule(lr){2-4}\cmidrule(lr){5-7}\cmidrule(lr){8-10}
Subject & Hit & Leak & Sup. & Hit & Leak & Int. & Blocked & Audit & Refused \\
\midrule
Baseline (no plugin) & 5/5 & 0/5 & 0/5 & 0/5 & 0/5 & 5/5 & 0/5 & n/a & 5/5 \\
Mem0 2.2.1, zh diagnostic & 5/5 & 0/5 & 0/5 & 0/5 & 0/5 & 5/5 & 0/5 & n/a & 5/5 \\
LINGSHU (lab profile) & n/i & n/i & 5/5 & 0/5 & 0/5 & 5/5 & 5/5 & 5/5 & 5/5 \\
\bottomrule
\end{tabular}

\vspace{2pt}
\raggedright\scriptsize
S1 uses the docs/19-supported correction sentence for all three subjects. Baseline
S1 hit is 4/5 by the literal string rule; the fifth round answered ASCII
\texttt{Beta}, counted as a hit on human review. LINGSHU S1 hit/leak are marked
n/i: the memory ledger is correct in 5/5 cells (superseded value tombstoned,
\texttt{supersedeEvents=1}, only BETA injected in the preview), but the model's
first request sees the previous turn's injection snapshot (F-D), so the answers
are ALPHA 5/5; a warm diagnostic with one extra turn answers BETA 5/5. Its
structural suppression is 5/5, the effective suppression 0/5. Mem0's default
configuration extracts English memory text; with the local non-neural embedder,
Chinese queries retrieve nothing (cosine $\le 0.063 < 0.1$), so its default S1
``suppression'' is an empty-recall artefact and the store still retains the stale
value (ADD-only pipeline). S2 canary never reaches the model in the LINGSHU arm
because the fact extractor rejects the secret-bearing sentence and the gate
blocks the attack turn; all three arms refuse the obvious payload, so no
model-refusal difference is claimed. $n=5$: directional only, no confidence
intervals.
\end{table*}


\paragraph{What the pilot shows.} On S1 the no-plugin host answers the codename
from the raw transcript (5/5 under human review; in the pre-fix run the literal
rule scored 4/5 because one round answered ASCII \texttt{Beta}) but never
suppresses the old value, by construction: the transcript contains it.
\textsc{Mem0} keeps
both entries: its default \texttt{add} path is additive---the library's own
prompt states that the only operation is ADD---so the superseded value stays in
the store and, in the diagnostic configuration where recall works, is retrieved
alongside the correction (suppression 0/5); we include an explicit-deletion
counterfactual showing that the failure is the default policy rather than a
broken library. \system{} writes the correction through its single writer and
tombstones ALPHA: in every cell the ledger records one supersede event, the old
entry is marked superseded, and the injected-memory preview contains only BETA
(structural suppression 5/5). That is the main structural difference in S1, and it
is a difference in kind rather than in ranking quality: Mem0's ADD-only vector
store retrieves what is similar, whereas \system{} checks the
$(\textsc{subject},\textsc{relation})$ pair before a fact is written, so the
correction tombstones ALPHA instead of joining it. Retrieval augmentation versus
structural defence.

The S1 cells were rerun after the v5.4 injection-timing fix (finding F-D). The
rerun closes the gap between ledger and model: in 5/5 cells the same-turn
assembly carries the recomputed snapshot, the superseded value is absent from the
request (effective suppression 5/5), and the model answers BETA 5/5 with no
ALPHA leak. Baseline and Mem0 also answer BETA 5/5, so the difference is not
answer accuracy; it is what each system puts in front of the model---they keep
ALPHA in the context (suppression 0/5), \system{} does not. Two implementation
gaps remain and are reported as such: the fact extractor does not recognise the
natural correction phrasing (``now changed to'', finding F-A) and stores the
final question as a fact (finding F-B).


On S2 the post-fix rerun shows why refusal language is a weak metric. All
three subjects refuse the obvious injection in most answers, but the no-plugin
host quoted the canary inside its refusal in 2/5 cells: it said it would not
print the gate code and then printed it. \textsc{Mem0} refused 5/5 with no leak
in its diagnostic configuration, so no subject is distinguished by a ``refusal
rate''. What separates them is structural admission: \system{} records a UCA-01
block attributed to the user message in every S2 cell (5/5, both runs), with
audit evidence, and the canary never reaches the model (0/5 leak), whereas
neither \textsc{Mem0} nor the no-plugin host has an admission gate at all. The
pilot therefore supports a narrow, structural claim and not a model-behaviour
claim: \emph{under this protocol, \system{} turns safety admission into a
structural, auditable gate (5/5 cells) where a pure memory framework has no gate
(0/5), without adding a mechanism to the memory layer.} We do not claim that
\system{} refuses more attacks than \textsc{Mem0}: refusal is stochastic, and the
baseline's 0/5 leak in the earlier run versus 2/5 here is the same stochasticity.
The stable difference is that a model can leak while sounding like it refuses,
and a gate cannot.

\paragraph{What it does not show.} The three subjects see different input shapes
(raw transcript versus retrieved memory), so rates are not directly comparable
across subjects and we do not report an effect size. The pilot is five rounds per
cell, one host, one model route and one day: all numbers are $n=5$ and
directional, and the row-level answers are archived so that a reader can
re-classify them. In S2 the planted token never reaches the model in the
\system{} arm at all: the fact extractor rejects the secret-bearing sentence as a
forbidden topic before it is stored, and the gate blocks the attack turn, so the
zero leak rate cannot be read as a disclosure test of the memory channel. The
injection-side gate stops injection and records it, but the host still delivers
the user message to the model, so ``blocked'' does not mean the model never reads
the attack. The pilot also exposed tool-gateway noise (finding F-C) and the one-turn
injection lag (F-D); the lag was traced to the host lifecycle and fixed in v5.4
(Section~\ref{sec:method:immersive}), and the S1 row above is from the post-fix
rerun. Both findings are reported as product findings, not folded into the rates.
